\subsection{FAMILIES OF CONVEX FUNCTIONS}

PROPOSITION 2. Let \(f_{i}\) ( \(1 \leqslant i \leqslant p\) ) be \(p\) convex functions on an interval \(\mathbf{I} \subset \mathbf{R}\) , and \(c_{i}\) ( \(1 \leqslant i \leqslant p\) ) be \(p\) arbitrary positive numbers; then the function \(f = \sum_{i=1}^{p} c_{i} f_{i}\) is convex on \(\mathbf{I}\) . Further, if for at least one index \(j\) the function \(f_{j}\) is strictly convex on \(\mathbf{I}\) , and \(c_{j} > 0\) , then \(f\) is strictly convex on \(\mathbf{I}\) .

This follows immediately by applying the inequality (1) (resp. (2)) to each of the \(f_{i}\) , multiplying the inequality for \(f_{i}\) by \(c_{i}\) , and then adding term-by-term.

PROPOSITION 3. Let \((f_{\alpha})\) be a family of convex functions on an interval \(\mathbf{I} \subset \mathbf{R}\) ; if the upper envelope \(g\) of this family is finite at every point of \(\mathbf{I}\) then \(g\) is convex on \(\mathbf{I}\) .

Indeed, the set of points \((x, y) \in \mathbf{R}^{2}\) lying above the graph of g is the intersection of the convex sets formed by the points lying above the graph of each of the functions \(f_{\alpha}\) ; so it is convex.

PROPOSITION 4. Let H be a set of convex functions on an interval \(I \subset R\) ; if F is a filter on H which converges pointwise on I to a finite real function \(f_{0}\) , then this function is convex on I.

To see this it suffices to pass to the limit along F in the inequality (1).
% semantic-fragments: legacy-input-seam
\relax{}
